\documentclass[11pt]{article}
\usepackage[margin=0.9in]{geometry}
\usepackage{amsmath,amssymb,amsthm,mathtools,booktabs,microtype,xcolor,hyperref}
\hypersetup{colorlinks=true,linkcolor=blue!55!black,urlcolor=blue!55!black,citecolor=blue!55!black}
\newtheorem{theorem}{Theorem}
\newtheorem{lemma}[theorem]{Lemma}
\newtheorem{corollary}[theorem]{Corollary}
\theoremstyle{remark}\newtheorem{remark}[theorem]{Remark}
\newcommand{\Rea}{\operatorname{Re}}
\newcommand{\pp}[1]{\left(#1\right)_+}
\newcommand{\ind}{\mathbf{1}}
\newcommand{\C}{\mathcal C}
\newcommand{\Z}{\mathcal Z}
\title{Local prime-power spectral detection\\\large One-sided Mellin thresholds and conductor isolation after GRH}
\author{Research note}
\date{26 November 2026}
\begin{document}
\maketitle
\begin{abstract}
For any fixed multiplicative width $R>1$ we construct a four-scale, compactly supported signed prime-power statistic. Its Mellin multiplier cancels the prime-number main term but is nowhere zero inside the critical strip. For any finite conjugation-invariant positively weighted family of Dirichlet characters, the abscissa of convergence of the Mellin integral of the positive part of this statistic is exactly the supremum of the real parts of the family's nontrivial zeros. Under GRH the normalized statistic is asymptotically uniformly almost periodic, and positive square-root-sized excursions occupy a positive proportion of every sufficiently long logarithmic interval. A divisor inversion produces an exact statistic for primitive characters of one conductor, using only congruences and prime powers. These results are classical proofs using the reported GRH theorem only for the final specialization; they do not reprove GRH. The precise synthesis is an independently developed candidate for novelty, not a priority claim.
\end{abstract}

\section{Source audit and motivation}
The repository \cite{sernau} was inspected using its GitHub source at commit
\texttt{72fead1} (full hash in the bibliography). The root README describes two editions, with Lean v4.35.0 and Mathlib pinned to \texttt{56a00585} (full hash in the archive README). The readable guide identifies three main ingredients: short integral forms, enough forms via a second moment, and long integral forms via determinant estimates. Their combination yields the main inequality, then a one-sided weighted prime error bound. The following is the portion relevant here (paths relative to \texttt{grh/}).

\begin{center}\small
\begin{tabular}{p{.41\linewidth}p{.51\linewidth}}
\toprule
Source & Role in this note\\\midrule
\texttt{src/Statements.lean}, \texttt{src/Main.lean} & Mathlib-level RH and Dirichlet GRH statements; half-plane nonvanishing and critical-strip zeros.\\
\texttt{src/New/Reduction/Zeros/}\newline\texttt{PrimeError.lean} & Weighted ramp, character orthogonality, meromorphic transform with zero-order residues.\\
\texttt{src/New/Reduction/Zeros/}\newline\texttt{Criterion.lean} & One-sided bound implies zero exclusion, with a quadratic twist selected to remove competing residues.\\
\texttt{src/New/Reduction/PrimeSums.lean} & Exact comparison of prime sums at different scales.\\
\texttt{src/New/Reduction/OneSidedBound.lean} & Main inequality implies the analytic bound via auxiliary scales and errors.\\
\texttt{scripts/PrintAxioms.lean} & Prints axiom dependencies of the five exported theorems.\\\bottomrule
\end{tabular}
\end{center}

More specifically, the proof uses the remainder
\[
 R_{d,q}(x)=\sum_{\substack{n\le x\\n\equiv1\bmod q}}\Lambda(n)
 (1-\chi_d(n))(1/n-1/x)-\frac{\log x}{\varphi(q)}-c_{d,q}.
\]
Its transform in the half-plane $\Rea s>0$ has residues given by differences of zero orders divided by $\varphi(q)s(1-s)$. The criterion chooses an auxiliary prime $d$ so that twisted zeros do not cancel the residue at issue. The file \texttt{OneSidedBound.lean} bounds the remainder by a square-root-scale error with an arbitrary small exponent loss. Our transfer is of the \emph{analytic} ideas (ramp weights, scale comparison, zero-order residues, and one-sided Landau), not a claim that the automorphic determinant construction has been generalized.

The five exported statements are:\par\smallskip
{\small\noindent\texttt{riemannHypothesis}, \path{riemannZeta_eq_zero_re},\\
\path{LFunction_ne_zero_of_half_lt_re},\\
\path{LFunction_eq_zero_re}, \path{LFunction_eq_zero_re_of_isPrimitive}.\par\smallskip}
\noindent The last three concern Dirichlet characters, including imprimitive ones, not all automorphic $L$-functions. The README reports machine checking only, not a human reading of the proof. Rubinstein's independent verification page documents checks of the original proof and links to the shortened-proof check \cite{rubinstein}. Wikipedia's sections on the proof and prime distribution \cite{wiki} describe the same zero/error relationship. We did not rebuild the Lean repository or run a kernel checker. The new theorems below have conventional mathematical proofs.

\section{A compact four-scale filter}
Fix $r>1$ (the width is $R=r^3$) and define
\begin{align*}
 P_r(z)&=(z-1)^2(rz-1)=-1+(r+2)z-(2r+1)z^2+rz^3,\\
 w_r(u)&=\sum_{j=0}^3 c_j\pp{u^{-1}-r^{-j}},\quad (c_0,c_1,c_2,c_3)=(-1,r+2,-2r-1,r).
\end{align*}
In this notation $v_+=\max(v,0)$. Direct cancellation gives
\begin{equation}\label{kernel}
 w_r(u)=\begin{cases}
 0,&0<u\le1,\\
 u^{-1}-1,&1<u\le r,\\
 2/r-(r+1)/u,&r<u\le r^2,\\
 r/u-r^{-2},&r^2<u\le r^3,\\
 0,&u>r^3.
 \end{cases}
\end{equation}
It is continuous, negative up to $u=r(r+1)/2$ inside its support, then positive. The multiplier is
\begin{equation}\label{W}
 W_r(s)=\int_0^\infty w_r(u)u^{s-1}\,du
       =\frac{P_r(r^{s-1})}{s(s-1)}.
\end{equation}
The apparent singularities at $0,1$ are removable. Indeed for $\Rea s>1$ each ramp integral is $r^{j(s-1)}/[s(s-1)]$, and the compact support yields an entire continuation. In particular $W_r(1)=0$ and $W_r(s)=O_{r,I}((1+|\Im s|)^{-2})$ for $\Rea s$ in any compact interval $I$. For $0<\Rea s<1$,
\begin{equation}\label{nonzero}
 r^{-1}<|r^{s-1}|<1,\qquad W_r(s)\ne0.
\end{equation}
For real $0<s<1$, $P_r(r^{s-1})>0$, and hence $W_r(s)<0$.

For any real sequence $a(n)$ with $|a(n)|$ bounded, put
\[
 S_a(x)=\sum_{n\le x}\Lambda(n)a(n)(1/n-1/x),\quad
 G_{a,r}(x)=\sum_{n\ge1}\Lambda(n)a(n)w_r(n/x).
\]
(The von Mangoldt function is $\Lambda(p^k)=\log p$ and zero otherwise.) Exactly,
\begin{equation}\label{dilation}
 G_{a,r}(x)=x\sum_{j=0}^3 c_j S_a(r^j x).
\end{equation}
The sum needs only $x<n<r^3x$. All terms $A\log x+B+C/x$ in $S_a(x)$ are annihilated. No analytic constant is needed. In the dyadic case $r=2$, the four coefficients are $(-1,4,-5,2)$ and the middle sign change is at $3x$. Any prescribed fixed width $1+\delta>1$ is allowed by taking $r=(1+\delta)^{1/3}$; constants below need not be uniform as $r\downarrow1$.

\begin{lemma}[Minimality in the ramp ansatz]\label{min}
Among linear combinations of $S_a(t_jx)$ at distinct fixed positive scales, four is the smallest possible number of nonzero coefficients if all past terms are to cancel for arbitrary $a$ and the logarithmic main term is to vanish.
\end{lemma}
\begin{proof}
Cancellation of the ramp coefficients below all cutoffs requires $\sum c_j=\sum c_j/t_j=0$ (the constant and $1/n$ pieces are independent). Killing the logarithmic main term additionally requires $\sum c_j\log t_j=0$. For three distinct scales the points $(\log t_j,1/t_j)$ are not collinear by strict convexity of $e^{-t}$. Thus the three columns $(1,\log t_j,1/t_j)$ are linearly independent. For two or one scales the same follows from distinctness. Four scales suffice by the polynomial construction. This is a minimality statement only within this linear ramp ansatz.
\end{proof}

\section{The positive-part threshold}
Let $\C$ be a nonempty finite family of Dirichlet characters (possibly of different moduli), with strictly positive real weights $b_\chi$, closed under complex conjugation with $b_{\bar\chi}=b_\chi$. Include multiplicity by adding weights if a character occurs more than once. Let $\chi^*$ be the primitive character inducing $\chi$. Let $\Z_\C$ be the set of distinct nontrivial zeros of the primitive functions $L(s,\chi^*)$, and
\[
 m_\C(\rho)=\sum_{\chi\in\C} b_\chi\operatorname{ord}_{\rho}L(s,\chi^*),\qquad
 B_\C=\sup_{\rho\in\Z_\C}\Rea\rho.
\]
Nontrivial means $0<\Rea\rho<1$. Standard zero counting and the functional equation give $1/2\le B_\C\le1$; for imprimitive functions the additional Euler-factor zeros on $\Rea s=0$ are not included. Write
\[
 a(n)=\sum_{\chi\in\C} b_\chi\chi(n)\in\mathbb R,\qquad G(x)=G_{a,r}(x).
\]

\begin{theorem}[Exact one-sided Mellin threshold]\label{main}
For any $r>1$ and family as above,
\begin{align}
 G(x)&=\sum_{\rho\in\Z_\C} m_\C(\rho)
 \frac{P_r(r^{\rho-1})}{\rho(1-\rho)}x^\rho+O_{\C,r}(1),\qquad x\ge1,\label{ef}\\
 \inf\left\{\sigma\in\mathbb R:\int_1^\infty G(x)_+\frac{dx}{x^{\sigma+1}}<\infty\right\}&=B_\C.\label{threshold}
\end{align}
The series in \eqref{ef} converges absolutely for each $x$ and is locally uniform. Consequently
\begin{equation}\label{limsup}
 \limsup_{x\to\infty}\frac{\log(1+G(x)_+)}{\log x}=B_\C.
\end{equation}
In particular GRH for the family is equivalent to the finiteness of the positive-part integral in \eqref{threshold} for every $\sigma>1/2$, and to the one-sided bound $G(x)\le C\sqrt x$ for some $C$ and all sufficiently large $x$.
\end{theorem}

\begin{proof}[Proof of the explicit formula]
By the Euler product, orthogonality is not required for the following identity. For $\Rea s>1$,
\begin{align}
 F(s):=\int_0^\infty G(x)x^{-s-1}\,dx
 &=W_r(s)\sum_{n\ge1}\frac{\Lambda(n)a(n)}{n^s}\notag\\
 &=-W_r(s)\sum_{\chi\in\C}b_\chi\frac{L'}{L}(s,\chi).\label{F}
\end{align}
For the interchange substitute $u=n/x$ in the integral of each term; compact support and absolute convergence of $\sum \Lambda(n)n^{-\Rea s}$ suffice. The right side gives meromorphic continuation to $\Rea s>0$; it is regular at $s=1$, since $W_r(1)=0$. For each nontrivial zero $\rho$ it has the nonzero residue
\begin{equation}\label{res}
 m_\C(\rho)\frac{P_r(r^{\rho-1})}{\rho(1-\rho)}.
\end{equation}
The extra Euler factors of imprimitive functions do not vanish in the open right half-plane. In particular, every real pole in that half-plane is simple with strictly positive residue, and all nontrivial zeros give poles. There is no cancellation by character weights.

For completeness, the contour giving \eqref{ef} is as follows. First replace each $\chi$ by $\chi^*$; the difference in $G(x)$ is $O_{\C,r}(1)$ because the only affected $n$ are powers of the finitely many primes dividing the moduli. In $(x,r^3x)$ the number of powers of a fixed $p$ is at most $1+3\log r/\log p$, and $w_r$ is bounded. For a primitive character $\psi$, Mellin inversion of the continuous compactly supported $w_r$ on $\Rea s=2$ gives
\[
 \sum_n\Lambda(n)\psi(n)w_r(n/x)=\frac{1}{2\pi i}\int_{2-i\infty}^{2+i\infty} -W_r(s)\frac{L'}{L}(s,\psi)x^s\,ds.
\]
Use the classical Dirichlet $L$-function facts: the count of zeros with ordinates in a unit interval at height $T$ is $O_\psi(\log(2+|T|))$, and away from the zeros by $\gg_\psi1/\log T$ the logarithmic derivative in the fixed strip $-1/2\le\Rea s\le2$ is $O_\psi(\log^2 T)$ at large height (see \cite{davenport,mv,kou}). These follow from the standard local partial-fraction estimate for $L'/L$ and zero counting. On the line $\Rea s=-1/2$ the functional equation and the absolutely convergent logarithmic derivative at $1-s$ give $O_\psi(\log(2+|\Im s|))$. This line contains no zeros or poles. Choose upper and lower heights tending to infinity avoiding the zeros by the stated distance: in each interval of length one, the excluded intervals have total length less than one if the implied distance constant is sufficiently small. This can be done simultaneously for the finite family.

Shift the contour to $\Rea s=-1/2$. The horizontal integrals vanish by the $O(T^{-2})$ decay of $W_r$ and the logarithmic derivative bound, for fixed $x$. The remaining vertical integral is absolutely convergent and $O_{\psi,r}(x^{-1/2})$. The pole at $1$ is removed by $W_r(1)=0$; a possible trivial zero at $0$ contributes just a constant because $W_r$ is entire. All other enclosed poles are precisely nontrivial zeros, with residues \eqref{res} times $x^\rho$. The zero-count estimate makes the sum of absolute coefficients finite: at high ordinates each coefficient is $O_r(|\Im\rho|^{-2})$ per zero, and there are only finitely many small ordinates. Local uniform convergence follows as $|x^\rho|\le\max(1,x)$ on compact positive intervals. Summing over the finite family and restoring imprimitive characters proves \eqref{ef}.
\end{proof}

\begin{lemma}[Landau for a nonnegative Mellin density]\label{landau}
Let $h\ge0$ be locally integrable on $[1,\infty)$ and suppose $H(s)=\int_1^\infty h(x)x^{-s-1}\,dx$ has a finite real abscissa $\alpha$ of convergence. Then $H$ cannot extend holomorphically to a neighborhood of $\alpha$.
\end{lemma}
\begin{proof}
Set $t=\log x$, so $H(s)=\int_0^\infty h(e^t)e^{-st}\,dt$. If $H$ were holomorphic near $\alpha$, choose a small $\epsilon>0$ such that a disc of radius $2\epsilon$ about $s_0=\alpha+\epsilon$ lies in the union of that neighborhood and the convergence half-plane. The Taylor coefficients of $H(s_0-z)$ are $\int_0^\infty t^k h(e^t)e^{-s_0t}dt/k!$, all nonnegative (differentiation under the integral is justified strictly to the right of $\alpha$). The series converges at $z=3\epsilon/2$. By monotone convergence its value is the original Laplace integral at $\alpha-\epsilon/2$, a contradiction.
\end{proof}

\begin{proof}[Proof of the threshold and equivalences]
From \eqref{ef} we have $|G(x)|\ll_{\C,r}x^{B_\C}$, hence absolute convergence for $\sigma>B_\C$. Suppose instead the positive-part integral converges at some $\sigma<B_\C$. Choose $\theta$ with $\max(0,\sigma)<\theta<B_\C$. Put
\[
 H_\pm(s)=\int_1^\infty (\pm G(x))_+ x^{-s-1}\,dx.
\]
Then $H_+$ is holomorphic in $\Rea s>\theta$. The function $E(s)=\int_0^1G(x)x^{-s-1}dx$ is entire, since $G(x)=0$ for $x\le2/r^3$ and $G$ is continuous. Initially in $\Rea s>1$,
\[
 H_-(s)=H_+(s)-F(s)+E(s).
\]
If the abscissa $\alpha$ of $H_-$ exceeded $\theta$, it would be finite by the upper bound. The right side is a meromorphic continuation near $\alpha$. If regular there, it contradicts Lemma~\ref{landau}. Otherwise $F$ has a simple pole with positive real residue, by \eqref{res}; the continued expression for $H_-$ then tends to $-\infty$ along real $s\downarrow\alpha$ from above. This contradicts the nonnegativity of its convergent integral. (The identities extend by uniqueness of meromorphic continuation throughout the intersection of their half-planes.) Thus $H_-$ converges in $\Rea s>\theta$, even if its abscissa is $-\infty$. It follows that $F=H_+-H_-+E$ is holomorphic in that half-plane. But by the definition of $B_\C$ some $\rho$ has $\Rea\rho>\theta$, and \eqref{res} is nonzero. This contradiction proves \eqref{threshold}.

The upper bound gives the upper inequality in \eqref{limsup}. If the limsup were less than $B_\C$, choose $0<\eta<\sigma<B_\C$ above it. Then eventually $G(x)_+\le x^\eta$, implying convergence at $\sigma$, impossible. This proves equality. The functional equation for each primitive character reflects zeros by $\rho\mapsto1-\bar\rho$ in the same $L$-function, so $B_\C=1/2$ exactly when its zeros all lie on the line. The claimed equivalences follow from the threshold and the bound obtained from \eqref{ef}. No assumption excluding real exceptional zeros has been used; their positive residues are precisely what makes the negative divergence argument work.
\end{proof}

\section{Critical endpoint: positive density and spectral recovery}
Assume GRH for the family. Group distinct zeros by ordinates $\gamma$, including $\gamma=0$ if it occurs, and define
\begin{equation}\label{A}
 A_\gamma=\frac{m_\C(1/2+i\gamma)P_r(r^{-1/2+i\gamma})}{1/4+\gamma^2},\qquad M=\sum_\gamma|A_\gamma|.
\end{equation}
Set $A_\gamma=0$ if the ordinate does not occur. Then $A_{-\gamma}=\bar A_\gamma$, $A_0\ge0$, $0<M<\infty$, and each actual ordinate has $A_\gamma\ne0$. Every primitive Dirichlet $L$-function has infinitely many nontrivial zeros, by its standard zero-count asymptotic, so some nonzero $\gamma_*$ occurs. Put $f(t)=e^{-t/2}G(e^t)$ for $t\ge0$.

\begin{theorem}[Uniform positive excursions]\label{density}
Let $b=|A_{\gamma_*}|>0$ and $\varepsilon=b/4$. Then
\begin{align}
 \liminf_{H\to\infty}\inf_{U\ge0}\frac1H\int_U^{U+H}f(t)_+\,dt&\ge\frac{A_0+b}{2},\label{meanpositive}\\
 \liminf_{H\to\infty}\inf_{U\ge0}\frac1H\left|\{t\in[U,U+H]:f(t)>\varepsilon\}\right|&\ge\frac{b}{4M}>0.\label{densityformula}
\end{align}
Here $|\cdot|$ denotes Lebesgue measure. In particular the integral in \eqref{threshold} diverges at $\sigma=1/2$. More strongly, square-root-sized positive excursions have uniformly positive lower logarithmic density in all long logarithmic intervals. Also, for every real $v$,
\begin{equation}\label{recover}
 \lim_{H\to\infty}\frac1H\int_U^{U+H}f(t)e^{-ivt}\,dt=A_v,
\end{equation}
uniformly in $U\ge0$. Thus the local prime statistic determines every zero ordinate and its total weighted multiplicity without a linear-independence hypothesis on ordinates.
\end{theorem}
\begin{proof}
Formula \eqref{ef} gives $f(t)=g(t)+O(e^{-t/2})$, where $g(t)=\sum_\gamma A_\gamma e^{i\gamma t}$ is a uniformly absolutely convergent trigonometric series, bounded by $M$. For each fixed frequency $v$, the mean of any nonmatching exponential over $[U,U+H]$ has modulus at most $2/(H|\gamma-v|)$; the bound is independent of $U$. Approximation by a finite subseries proves \eqref{recover}, because the remainder has sup norm bounded by the coefficient tail. The mean of $O(e^{-t/2})$ is $O(1/H)$ uniformly in $U\ge0$.

Taking $v=0$ gives mean $A_0$. Taking $v=\gamma_*$ and using the triangle inequality shows the lower limit of the means of $|f|$ is at least $b$, uniformly in $U$. Since $2f_+=|f|+f$, this proves \eqref{meanpositive}. Pointwise $f_+(t)\le \varepsilon+M\ind_{\{f(t)>\varepsilon\}}+O(e^{-t/2})$ (choose a nonnegative absolute error bound). Integrating, taking the infimum and then lower limit proves \eqref{densityformula}. With $x=e^t$, the critical integral is $\int_0^\infty f(t)_+dt$, which diverges by \eqref{meanpositive} at $U=0$. Recovery of multiplicities uses the nonvanishing multiplier \eqref{nonzero} and the known factor $1/4+v^2$. This is a uniqueness theorem for the infinite exact arithmetic signal, not a finite numerical certification of zeros.
\end{proof}

\section{Arithmetic specializations: one class, one conductor}
\begin{corollary}[Identity residue class]\label{class}
For any $q\ge1$, let
\[
 G_{q,r}(x)=\sum_{\substack{x<n<r^3x\\n\equiv1\bmod q}}\Lambda(n)w_r(n/x).
\]
Theorem~\ref{main} holds with $B_\C=B_q$, the supremum over all Dirichlet characters mod $q$, since $a(n)=\varphi(q)^{-1}\sum_{\chi\bmod q}\chi(n)$ is the indicated residue class. The reported Dirichlet GRH implies $|G_{q,r}(x)|\ll_{q,r}\sqrt x$ and Theorem~\ref{density} for every fixed $q,r$. Conversely even just $\int_1^\infty (G_{q,r}(x))_+x^{-\sigma-1}dx<\infty$ for every $\sigma>1/2$ forces GRH for every character mod $q$. At $q=1$ this is RH.
\end{corollary}
\begin{proof} Character orthogonality, followed by Theorems~\ref{main} and \ref{density}. \end{proof}

\begin{corollary}[Primitive-conductor isolation]\label{conductor}
Define the integer arithmetic weight
\begin{equation}\label{primweight}
 a_q^{\rm prim}(n)=\ind_{(n,q)=1}\sum_{d\mid q}\mu(q/d)\varphi(d)\ind_{n\equiv1\bmod d},
\end{equation}
and the local statistic $J_{q,r}(x)=\sum_n\Lambda(n)a_q^{\rm prim}(n)w_r(n/x)$. If at least one primitive character mod $q$ exists, Theorems~\ref{main} and \ref{density} hold for $J_{q,r}$ with $B_\C$ equal to the supremum over \emph{only} the primitive characters of conductor $q$. In particular the exact threshold is $1/2$ under Dirichlet GRH, and the positive-density and multiplicity-recovery conclusions hold. If there are no primitive characters, $J_{q,r}\equiv0$ instead.
\end{corollary}
\begin{proof}
For $(n,q)=1$, every character mod $d\mid q$ is induced by exactly one primitive character mod $e\mid d$, with the same value at $n$. Orthogonality yields
\[
 \varphi(d)\ind_{n\equiv1\bmod d}=\sum_{e\mid d}\sum_{\substack{\chi\bmod e\\\chi\text{ primitive}}}\chi(n).
\]
Divisor M\"obius inversion in $d$ proves that the inner sum in \eqref{primweight} equals the sum of primitive characters mod $q$. When $(n,q)>1$ those characters all vanish, as does \eqref{primweight}. Thus the identity holds on every integer, including prime powers at primes dividing $q$. The family of primitive characters is conjugation-stable, all weights are $1$, and the theorems apply directly. The empty case follows from the same identity.
\end{proof}

\paragraph{Example.} At $q=5$, $a_5^{\rm prim}(n)=3$ for $n\equiv1\pmod5$, $-1$ for $n\equiv2,3,4\pmod5$, and $0$ for $5\mid n$. This statistic removes the zeta contribution exactly while retaining all three primitive characters mod $5$. For $q=6$ there are no primitive characters and the weight vanishes identically. The coprimality indicator in \eqref{primweight} is essential for exactness.

\paragraph{Why the conductor result is stronger than an error bound.} For an individual $q$, all lower-conductor spectral lines are absent unless they happen to coincide with actual conductor-$q$ lines. The transform residue is a sum of positive multiplicities at coincident zeros. This assertion does not require simplicity or distinctness of zeros across characters. Using \eqref{recover} on $J_{q,r}$ gives the aggregate multiplicity at conductor $q$, but not which primitive character has a given zero.

\section{Verification, originality limits, and research directions}
The accompanying Python script checks the piecewise ramp identity and cancellation exactly using rational arithmetic for several ratios, verifies the primitive-conductor weights by an independent character enumeration for small moduli, and produces sample local sums with floating-point logarithms. These are finite algebraic/computational checks, not an analytic proof or a Lean certification. The mathematical proof is Theorems~\ref{main}--\ref{density} and the corollaries, using classical Dirichlet $L$-function theory as cited.

\paragraph{Prior art and scope.} Landau's sign principle, the explicit formula, character orthogonality and M\"obius inversion are classical. Suzuki \cite{suzuki} establishes sign criteria for RH and Dirichlet GRH using other weighted Mangoldt sums. The reported Alp\"oge proof also uses positive families and one-sided/windowed criteria; the condensed repository's twist construction is not needed for our chosen arithmetic statistic. Rubinstein's explanatory draft \cite{rubdraft} discusses a centered ramp and one-sided RH. A closely related four-scale finite prime-power window using a different ramp and an integer LCM ratio is described by Dumas \cite{dumas}; therefore the bare idea of a compact cubic filter is not claimed as new. Our candidate contribution is the combined exact positive-part Mellin threshold, arbitrary fixed width, conductor-isolating formula, and uniform critical-endpoint density and spectral recovery. A targeted literature search did not establish priority for this precise combination. No claim is made that the work resolves a separate famous open conjecture, or that all of its ingredients are original.

\paragraph{Further research questions.} Can one optimize the ratio and kernel for a specified conductor and provable lower density? Can a uniform-in-$q$ and shrinking-width variant be established with explicit constants? Can bounded-error finite data recover multiplicity robustly when frequencies are close? Can the Landau proof be lifted into the repository's \texttt{IsPolar} API, and its explicit formula developed in Mathlib? Can the same compact filter be used for other $L$-functions with meromorphic logarithmic derivatives and positive zero multiplicities? Such an extension would need new analytic inputs and its own zero-line theorem; the Dirichlet GRH result alone is not sufficient.

\begin{thebibliography}{9}\small
\bibitem{sernau} L. Sernau et al., \emph{GRH for Dirichlet L-functions, in Lean 4}, 25 November 2026, pinned commit \texttt{72fead1}. \url{https://github.com/lsernau/grh-lean/tree/72fead102f663d5c7a407964a13dc0cf9d977207}.
\bibitem{rubinstein} M. Rubinstein, \emph{Alp\"oge/Claude proof of the Generalized Riemann Hypothesis}, verification and discussion, accessed 26 November 2026. \url{https://www.math.uwaterloo.ca/~mrubinst/alpoge_claude_grh/}.
\bibitem{wiki} Wikipedia, \emph{Riemann hypothesis}, sections on formal proof, consequences and prime distribution, accessed 26 November 2026. \url{https://en.wikipedia.org/wiki/Riemann_hypothesis}.
\bibitem{davenport} H. Davenport, \emph{Multiplicative Number Theory}, 3rd ed., revised by H. L. Montgomery, Springer, 2000. Chapters on Dirichlet $L$-functions and explicit formulae.
\bibitem{mv} H. L. Montgomery and R. C. Vaughan, \emph{Multiplicative Number Theory I: Classical Theory}, Cambridge University Press, 2007. Especially \S15.1 for Landau's theorem.
\bibitem{kou} D. Koukoulopoulos, \emph{The Distribution of Prime Numbers}, AMS, 2019, Chapter 11, in particular Lemma 11.4 (local zero count and logarithmic derivative). \url{https://dms.umontreal.ca/~koukoulo/documents/publications/primes.pdf}.
\bibitem{suzuki} M. Suzuki, On variants of Chebyshev's conjecture, \emph{Ramanujan J.} 68, 95 (2025), \url{https://doi.org/10.1007/s11139-025-01238-9}, arXiv:2411.07436.
\bibitem{rubdraft} M. Rubinstein, \emph{The Alp\"oge--Claude proof of the Riemann hypothesis explained for the Riemann zeta function}, working draft v2, 22 November 2026, linked from \cite{rubinstein}.
\bibitem{dumas} A. Dumas, \emph{In Pursuit of a Proof of the Riemann Hypothesis: From Schur minors and heat collisions to Gaussian detectors and signed primes}, public synopsis, equations (S11.15)--(S11.17), accessed 26 November 2026. \url{https://www.academia.edu/181655129/In_Pursuit_of_a_Proof_of_the_Riemann_Hypothesis_From_Schur_minors_and_heat_collisions_to_Gaussian_detectors_and_signed_primes}. Only the publicly indexed excerpt was inspected; its general claims have not been audited here.
\end{thebibliography}
\end{document}
